\section*{Chapter \ref{ch:rc:pnl-explain-and-independent-price-verification} --- P\&L Explain and Independent Price Verification}

\begin{solution}{exo:rc:pnl-explain-and-independent-price-verification:1}
$1\,086\times30\times12 = \text{USD}~391\,087$ (with the unrounded vanna).
\end{solution}

\begin{solution}{exo:rc:pnl-explain-and-independent-price-verification:2}
The strikes are symmetric around the forward and the model is normal: the long payer and the short
receiver have equal and opposite gamma, vega and theta, and equal delta of the same sign.
\end{solution}

\begin{solution}{exo:rc:pnl-explain-and-independent-price-verification:3}
The band is 33 to 37; the mark moves to the boundary, 37, an adjustment of $-1$.
\end{solution}

\begin{solution}{exo:rc:pnl-explain-and-independent-price-verification:4}
Each step is revalued at the factors already moved, so cross effects are credited to whichever factor
moves later. Averaging over all orders (the Shapley allocation) removes the dependence at the cost of
more revaluations; reporting the cross effect separately is the practical alternative.
\end{solution}

\begin{solution}{exo:rc:pnl-explain-and-independent-price-verification:5}
The cross term is proportional to the rate move: about half, USD~181 thousand for a 15 basis point move.
\end{solution}

\begin{solution}{exo:rc:pnl-explain-and-independent-price-verification:6}
Its fair value nets the long and short legs (zero here), but the uncertainty of each leg's price does not
net: each depends on a volatility whose consensus is dispersed. The AVA scales with the gross risk.
\end{solution}

\begin{solution}{exo:rc:pnl-explain-and-independent-price-verification:7}
The individual AVAs double (382\,470 and 468\,316) and the aggregated figure to USD~425\,393; the
simplified figure does not change.
\end{solution}

\begin{solution}{exo:rc:pnl-explain-and-independent-price-verification:8}
Thresholds allow marks anywhere inside the band; systematic marks at its favourable edge pass every
test, as in 2012. IPV also covers only instruments with independent prices; level~3 positions and
model inputs need other checks, and reserves and AVAs are still required.
\end{solution}

\begin{solution}{pb:rc:pnl-explain-and-independent-price-verification:1}
\textbf{1.} USD~4\,571\,501 actual; 4\,173\,095 explained.
\textbf{2.} USD~398\,406, 9\% of the P\&L.
\textbf{3.} USD~139\,103 per basis point; 1\,086 per basis point of rate and of volatility.
\textbf{4.} USD~391\,087; a residual of 7\,319.
\textbf{5.} Third-order terms: the change of vanna with the rate and the volatility over such a large move.
\textbf{6.} Booking and data (missing trades, stale market data), then missing risks (cross terms,
unmarked factors), then marks (IPV breaks); the first two are quickest to rule in or out.
\textbf{7.} Moving the factors one at a time shows that the second factor moved carries an effect that
depends on the first: a cross effect.
\textbf{8.} The volatility step contributes nothing, because the vega is zero at the start; all the cross
effect shows up in the rate step.
\textbf{9.} The cross term is a product of two moves; on quiet days both are small and so is their product.
\textbf{10.} A flat mark implies no skew; with the market's skew, a rise in rates changes the volatility at
each strike, and the explain would carry the effect as vega; without it, it appears as vanna that the desk
does not report.
\textbf{11.} The receiver's volatility is moved from 95 to 93 basis points: the book's value rises by
USD~146\,482.
\textbf{12.} USD~211\,771 under the core approach (190\,578 and 232\,964 aggregated at 50\%); 6\,872 under the
simplified one.
\textbf{13.} A money threshold tied to the desk's typical daily P\&L and a share (say a few per cent), with
automatic investigation on big-move days.
\textbf{14.} No: revaluation adds up by construction and hides missing risks; the risk-based explain tests
whether the desk's risk measures describe its P\&L.
\textbf{15.} The \omterm{def:rc:pnl-explain-and-independent-price-verification:unexplained}{unexplained P\&L}, its cause, the missing vanna in the desk's risk reports, and the fix
(skew marks and cross-Greek reporting).
\textbf{16.} When it does not scale with market moves, appears on quiet days, or reverses the next day: a
trade booked late, a wrong notional, a stale fixing.
\textbf{17.} Because the desk's pay depends on the P\&L it would explain and the marks it would verify.
\textbf{18.} The attribution test compares the risk model's P\&L with the pricing systems'; a desk whose
risk model lacks the cross term would show the same gap there.
\textbf{19.} Named result: \emph{the unexplained 400 thousand}: the missing vanna term, 1\,086 dollars per basis
point squared times 30 and 12, explains USD~391\,087 of the 398\,406 unexplained, leaving 7\,319.
\textbf{20.} Which risk, data or mark the desk's picture of itself is missing.
\end{solution}

\begin{solution}{iq:rc:pnl-explain-and-independent-price-verification:1}
Delta, gamma, vega (and volga), theta, the cross terms that matter (vanna, \omterm{def:rc:rates-risk:crossgamma}{cross-gamma} between curves),
carry and funding, new trades and amendments, reserves and fees; the remainder is unexplained.
\iqlookfor{Greeks, time, new business and the remainder.}
\end{solution}

\begin{solution}{iq:rc:pnl-explain-and-independent-price-verification:2}
Missing or stale market data, booking errors, missing \omterm{def:rc:market-risk-measures:factor}{risk factors} or cross terms, model changes,
large moves beyond the expansion, reserves and fees not included, and marks that changed without market
moves.
\iqlookfor{data, booking, model and marks.}
\end{solution}

\begin{solution}{iq:rc:pnl-explain-and-independent-price-verification:3}
Verifying marks and model inputs against independent sources and adjusting those outside tolerance. It is
hard for illiquid instruments and unobservable inputs, when consensus is thin or stale, and when thresholds
invite marks at their edges.
\iqlookfor{the process and its weak points.}
\end{solution}

\begin{solution}{iq:rc:pnl-explain-and-independent-price-verification:4}
Fair value is an exit price with market participants' assumptions; \omterm{def:rc:pnl-explain-and-independent-price-verification:pv}{prudent valuation} takes a conservative
point (90\% confidence) of the range of plausible exit values and deducts the difference from capital.
\iqlookfor{the confidence level and the capital deduction.}
\end{solution}

\begin{solution}{iq:rc:pnl-explain-and-independent-price-verification:5}
Start-of-day Greeks by trade and factor from the pricing library, end-of-day market data snapshots,
attribution by factor with thresholds and drill-down, revaluation fallback for large moves, alerts,
audit trail, and reconciliation with the general ledger.
\iqlookfor{data, computation, thresholds and controls.}
\end{solution}

\begin{solution}{iq:rc:pnl-explain-and-independent-price-verification:6}
Compare marks with consensus over time: the distribution of the mark's position within the band, its
correlation with the desk's P\&L targets, jumps at month end, and differences between month-end and
intramonth marks.
\iqlookfor{statistical tests on the position of marks.}
\end{solution}
